\documentclass[UTF8,11pt]{ctexart}
\usepackage[a4paper,margin=24mm]{geometry}
\usepackage{amsmath,amssymb,amsthm,mathtools,mathrsfs}
\usepackage[all]{xy}
\usepackage{xurl,hyperref,enumitem}
\usepackage{tikz}
\usetikzlibrary{arrows.meta,cd,decorations.markings,calc,patterns}
\hypersetup{hidelinks,breaklinks=true}
\setlength{\emergencystretch}{3em}
\allowdisplaybreaks
\newtheorem*{theorem}{Theorem}
\newtheorem*{thm}{Theorem}
\newtheorem*{lemma}{Lemma}
\newtheorem*{proposition}{Proposition}
\newtheorem*{prop}{Proposition}
\newtheorem*{corollary}{Corollary}
\newtheorem*{cor}{Corollary}
\newtheorem*{claim}{Claim}
\theoremstyle{definition}
\newtheorem*{definition}{Definition}
\newtheorem*{defn}{Definition}
\newtheorem*{remark}{Remark}
\newtheorem*{example}{Example}
\newcommand{\R}{\mathbb R}
\newcommand{\C}{\mathbb C}
\newcommand{\Z}{\mathbb Z}
\newcommand{\Q}{\mathbb Q}
\newcommand{\N}{\mathbb N}
\newcommand{\F}{\mathbb F}
\newcommand{\D}{\mathbb D}
\newcommand{\bB}{\mathbb B}
\newcommand{\bC}{\mathbb C}
\newcommand{\bR}{\mathbb R}
\newcommand{\bZ}{\mathbb Z}
\newcommand{\bN}{\mathbb N}
\newcommand{\bQ}{\mathbb Q}
\newcommand{\bD}{\mathbb D}
\newcommand{\bF}{\mathbb F}
\newcommand{\CP}{\mathbb{CP}}
\newcommand{\RP}{\mathbb{RP}}
\newcommand{\sO}{\mathscr O}
\newcommand{\sI}{\mathscr I}
\newcommand{\sF}{\mathscr F}
\newcommand{\sG}{\mathscr G}
\newcommand{\sH}{\mathscr H}
\newcommand{\sR}{\mathscr R}
\newcommand{\sM}{\mathscr M}
\newcommand{\sV}{\mathscr V}
\newcommand{\sU}{\mathscr U}
\DeclareMathOperator{\Hom}{Hom}
\DeclareMathOperator{\Ext}{Ext}
\DeclareMathOperator{\Tor}{Tor}
\DeclareMathOperator{\Spec}{Spec}
\DeclareMathOperator{\Proj}{Proj}
\DeclareMathOperator{\supp}{supp}
\DeclareMathOperator{\im}{im}
\DeclareMathOperator{\id}{id}
\DeclareMathOperator{\Tr}{Tr}
\DeclareMathOperator{\tr}{tr}
\DeclareMathOperator{\rank}{rank}
\DeclareMathOperator{\ord}{ord}
\DeclareMathOperator{\dist}{dist}
\DeclareMathOperator{\End}{End}
\DeclareMathOperator{\Aut}{Aut}
\DeclareMathOperator{\Gal}{Gal}
\DeclareMathOperator{\vol}{vol}
\DeclareMathOperator{\Cl}{Cl}
\DeclareMathOperator{\coker}{coker}
\DeclareMathOperator{\codim}{codim}
\DeclareMathOperator{\divergence}{div}
\DeclareMathOperator{\Ric}{Ric}
\DeclareMathOperator{\grad}{grad}
\DeclareMathOperator{\Hess}{Hess}
\newcommand{\abs}[1]{\left\lvert#1\right\rvert}
\newcommand{\sabs}[1]{\lvert#1\rvert}
\newcommand{\babs}[1]{\bigl\lvert#1\bigr\rvert}
\newcommand{\Babs}[1]{\Bigl\lvert#1\Bigr\rvert}
\newcommand{\bbabs}[1]{\biggl\lvert#1\biggr\rvert}
\newcommand{\BBabs}[1]{\Biggl\lvert#1\Biggr\rvert}
\newcommand{\norm}[1]{\left\lVert#1\right\rVert}
\newcommand{\snorm}[1]{\lVert#1\rVert}
\newcommand{\bnorm}[1]{\bigl\lVert#1\bigr\rVert}
\newcommand{\Bnorm}[1]{\Bigl\lVert#1\Bigr\rVert}
\newcommand{\bbnorm}[1]{\biggl\lVert#1\biggr\rVert}
\newcommand{\BBnorm}[1]{\Biggl\lVert#1\Biggr\rVert}
\newcommand{\widebar}[1]{\overline{#1}}
\newcommand{\nicefrac}[2]{\frac{#1}{#2}}
\newcommand{\linnprod}[2]{\langle#1,#2\rangle}
\newcommand{\myindex}[1]{#1}
\newcommand{\glsadd}[1]{}
\newcommand{\avoidbreak}{}
\newcommand{\sourceref}[1]{\textnormal{[原文：\nolinkurl{#1}]}}
\newcommand{\thmref}[1]{\sourceref{#1}}
\newcommand{\lemmaref}[1]{\sourceref{#1}}
\newcommand{\propref}[1]{\sourceref{#1}}
\newcommand{\corref}[1]{\sourceref{#1}}
\newcommand{\defnref}[1]{\sourceref{#1}}
\newcommand{\exampleref}[1]{\sourceref{#1}}
\newcommand{\exerciseref}[1]{\sourceref{#1}}
\newcommand{\figureref}[1]{\sourceref{#1}}
\newcommand{\sectionref}[1]{\sourceref{#1}}
\newcommand{\appendixref}[1]{\sourceref{#1}}
\newcommand{\stacksref}[2]{\href{https://stacks.math.columbia.edu/tag/#1}{#2 [#1]}}

\begin{document}
\section*{048\quad 两素因子群阶的Burnside可解性定理}
\subsection*{来源与计数目标}
Pavel Etingof 与 Oleg Golberg、Sebastian Hensel、Tiankai Liu、Alex Schwendner、Elena Yudovina、Dmitry Vaintrob 合著的 \emph{Introduction to Representation Theory} 课程讲义。MIT 18.712，Fall 2010 的 OCW 原讲义版本，不是限制转载的 2016 年 AMS 出版书。Chapter 4，\S4.5，Theorem 4.20 及 Theorems 4.21、4.23、Lemmas 4.22、4.24 的完整证明，分章 PDF 第6--8页；另收第5页 Proposition 4.17 的中心特征标计算。合订本相应为印刷页51--54。。获取日期：2026-09-28。
\par\url{https://ocw.mit.edu/courses/18-712-introduction-to-representation-theory-fall-2010/84358595a02a73bced2c4e363a5d66f0_MIT18_712F10_ch4.pdf}
\par\url{https://ocw.mit.edu/courses/18-712-introduction-to-representation-theory-fall-2010/pages/lecture-notes/}
\par\textbf{本文件只计一个问题：}$|G|=p^aq^b$ 的有限群可解。素数幂共轭类定理与代数整数引理作为同一证明链，不另计题数。
\subsection*{作者命题与人类证明}

\paragraph{Supporting central-character calculation (Proposition 4.17).}
Let $C$ be a conjugacy class in $G$, and $V$ an irreducible complex representation. The number $\lambda=|C|\chi_V(g)/\dim V$, for $g\in C$, is an algebraic integer.
\begin{proof}
Let $C$ be a conjugacy class in $G$, and $P=\sum_{h\in C}h$. Then $P$ is a central element of $\Z[G]$, so it acts on $V$ by some scalar $\lambda$, which is an algebraic integer (indeed, since $\Z[G]$ is a finitely generated $\Z$-module, any element of $\Z[G]$ is integral over $\Z$, i.e., satisfies a monic polynomial equation with integer coefficients). On the other hand, taking the trace of $P$ in $V$, we get $|C|\chi_V(g)=\lambda\dim V$, $g\in C$, so $\lambda=\frac{|C|\chi_V(g)}{\dim V}$.
\end{proof}
\begin{definition}[4.18]
A group $G$ is called solvable if there exists a series of nested normal subgroups
\[
\{e\}=G_1\triangleleft G_2\triangleleft\cdots\triangleleft G_n=G
\]
where $G_{i+1}/G_i$ is abelian for all $1\leq i\leq n-1$.
\end{definition}
\begin{theorem}[4.20 (Burnside)]
Any group $G$ of order $p^aq^b$, where $p$ and $q$ are prime and $a,b\geq0$, is solvable.
\end{theorem}
Before proving Burnside's theorem we will prove several other results which are of independent interest.
\begin{theorem}[4.21]
Let $V$ be an irreducible representation of a finite group $G$ and let $C$ be a conjugacy class of $G$ with $\gcd(|C|,\dim(V))=1$. Then for any $g\in C$, either $\chi_V(g)=0$ or $g$ acts as a scalar on $V$.
\end{theorem}
The proof will be based on the following lemma.
\begin{lemma}[4.22]
If $\varepsilon_1,\varepsilon_2,\ldots,\varepsilon_n$ are roots of unity such that $\frac1n(\varepsilon_1+\varepsilon_2+\cdots+\varepsilon_n)$ is an algebraic integer, then either $\varepsilon_1=\cdots=\varepsilon_n$ or $\varepsilon_1+\cdots+\varepsilon_n=0$.
\end{lemma}
\begin{proof}
Let $a=\frac1n(\varepsilon_1+\cdots+\varepsilon_n)$. If not all $\varepsilon_i$ are equal, then $|a|<1$. Moreover, since any algebraic conjugate of a root of unity is also a root of unity, $|a'|\leq1$ for any algebraic conjugate $a'$ of $a$. But the product of all algebraic conjugates of $a$ is an integer. Since it has absolute value $<1$, it must equal zero. Therefore, $a=0$.
\end{proof}
\begin{proof}[Proof of Theorem 4.21]
Let $\dim V=n$. Let $\varepsilon_1,\varepsilon_2,\ldots,\varepsilon_n$ be the eigenvalues of $\rho_V(g)$. They are roots of unity, so $\chi_V(g)$ is an algebraic integer. Also, by Proposition 4.17, $\frac1n|C|\chi_V(g)$ is an algebraic integer. Since $\gcd(n,|C|)=1$, there exist integers $a,b$ such that $a|C|+bn=1$. This implies that
\[
\frac{\chi_V(g)}n=\frac1n(\varepsilon_1+\cdots+\varepsilon_n)
\]
is an algebraic integer. Thus, by Lemma 4.22, we get that either $\varepsilon_1=\cdots=\varepsilon_n$ or $\varepsilon_1+\cdots+\varepsilon_n=\chi_V(g)=0$. In the first case, since $\rho_V(g)$ is diagonalizable, it must be scalar. In the second case, $\chi_V(g)=0$. The theorem is proved.
\end{proof}
\begin{theorem}[4.23]
Let $G$ be a finite group, and let $C$ be a conjugacy class in $G$ of order $p^k$ where $p$ is prime and $k>0$. Then $G$ has a proper nontrivial normal subgroup (i.e., $G$ is not simple).
\end{theorem}
\begin{proof}
Choose an element $g\in C$. Since $g\ne e$, by orthogonality of columns of the character table,
\[
\sum_{V\in\operatorname{Irr}G}\dim V\,\chi_V(g)=0.\tag{4}
\]
We can divide $\operatorname{Irr}G$ into three parts:
\begin{enumerate}
\item the trivial representation,
\item $D$, the set of irreducible representations whose dimension is divisible by $p$, and
\item $N$, the set of non-trivial irreducible representations whose dimension is not divisible by $p$.
\end{enumerate}
\begin{lemma}[4.24]
There exists $V\in N$ such that $\chi_V(g)\ne0$.
\end{lemma}
\begin{proof}
If $V\in D$, the number $\frac1p\dim(V)\chi_V(g)$ is an algebraic integer, so
\[
a=\frac1p\sum_{V\in D}\dim(V)\chi_V(g)
\]
is an algebraic integer. Now, by (4), we have
\[
\begin{aligned}
0&=\chi_{\C}(g)+\sum_{V\in D}\dim V\,\chi_V(g)+\sum_{V\in N}\dim V\,\chi_V(g)\\
 &=1+pa+\sum_{V\in N}\dim V\,\chi_V(g).
\end{aligned}
\]
This means that the last summand is nonzero.
\end{proof}
Now pick $V\in N$ such that $\chi_V(g)\ne0$; it exists by Lemma 4.24. Theorem 4.21 implies that $g$ (and hence any element of $C$) acts by a scalar in $V$. Now let $H$ be the subgroup of $G$ generated by elements $ab^{-1}$, $a,b\in C$. It is normal and acts trivially in $V$, so $H\ne G$, as $V$ is nontrivial. Also $H\ne1$, since $|C|>1$.
\end{proof}
\begin{proof}[Proof of Burnside's theorem]
Assume Burnside's theorem is false. Then there exists a nonsolvable group $G$ of order $p^aq^b$. Let $G$ be the smallest such group. Then $G$ is simple, and by Theorem 4.23, it cannot have a conjugacy class of order $p^k$ or $q^k$, $k\geq1$. So the order of any conjugacy class in $G$ is either $1$ or is divisible by $pq$. Adding the orders of conjugacy classes and equating the sum to $p^aq^b$, we see that there has to be more than one conjugacy class consisting just of one element. So $G$ has a nontrivial center, which gives a contradiction.
\end{proof}

\subsection*{整理说明（非作者正文）}
正文按作者支撑结论及证明次序保留；Proposition 4.17 的结论独立摘写出其原定义中的 $\lambda$，避免依赖被截掉的记号。未收与本题无关的 Frobenius 整除定理及其证明。标准先修：复表示完全可约与 Schur 引理、特征标列正交性、代数整数成环且有理代数整数为整数、代数共轭与范数、可解群对正规扩张封闭及类方程。独立工作是从共轭类大小互素性得到平均单位根的整性，再构造非平凡正规子群并排除最小反例，不能把 Burnside 定理作为先修。已取得正文文本，并查看分章 PDF 第6、8页图像；第5、7页截图返回错误，相关分式与非零符号另以同站合订本的同段文本交叉定位，未冒称全页图像检查。
\subsection*{可修改的简短标注（非作者正文）}
$c$：有限群的结构与可解性。

$a$：复特征标；共轭类和；代数整数；单位根平均与代数共轭；列正交性；表示核；最小反例与类方程。

\texttt{cross\_theory\_status = yes}。共轭类的整数大小经中心特征标转为代数整数约束，迫使某表示上的共轭类元素成为标量，再返回群内正规子群和可解性。位置：4.21--4.24 与 4.20 的末段证明。

全部标注均为非穷尽、可修改初稿；未标注不表示未使用。
\subsection*{许可与修改记录}
材料来自 MIT OpenCourseWare，作者与课程见来源。按该站所示 CC BY-NC-SA 4.0 许可复用，整理部分沿用同一许可。\url{https://creativecommons.org/licenses/by-nc-sa/4.0/}。2026-09-28：助手转录题证、调整数学排版并加入来源、范围和初稿标注；不暗示作者或 MIT 认可本次整理。所留为本次题证转录摘录，不是原 PDF 下载件。
\end{document}
